\section{模型假设}

为突出低温防护服传热与优化的主要机制，本文作如下假设：

\begin{itemize}[label=\textbullet,leftmargin=3em,labelsep=0.75em,itemsep=0.25ex,topsep=0.25ex]
\item \textbf{一维层合平板假设：}三层衣料厚度远小于人体表面特征尺度，各层有效传热面积近似相等，忽略面内温差，仅考虑沿厚度方向的传热。

\item \textbf{人体恒温边界假设：}暴露期间人体温度保持$37\,^{\circ}\mathrm C$，人体与内层织物之间的换热采用题设对流换热系数$h_b$描述，并以贴身侧温度作为受冷评价量。

\item \textbf{理想层间接触假设：}忽略各材料层之间的接触热阻，界面处温度连续且热流连续。

\item \textbf{辐射忽略假设：}依据题设仅考虑传导与对流换热，忽略人体、防护服与外界之间的辐射换热。
\end{itemize}